\documentclass[11pt]{article}
\usepackage[margin=1in]{geometry}
\usepackage{amsmath,amssymb}
\usepackage{xcolor}
\usepackage{booktabs}
\usepackage{enumitem}
\usepackage{fancyhdr}
\usepackage{tcolorbox}
\tcbuselibrary{skins,breakable}
\usepackage{hyperref}
\hypersetup{colorlinks=true, linkcolor=blue, urlcolor=blue}

\definecolor{primary}{RGB}{30,64,140}
\definecolor{tipbg}{RGB}{235,244,255}
\definecolor{examplebg}{RGB}{240,248,235}
\definecolor{formulabg}{RGB}{255,247,224}
\definecolor{warnbg}{RGB}{255,236,236}

\pagestyle{fancy}
\fancyhf{}
\renewcommand{\headrulewidth}{0.4pt}
\cfoot{\thepage}

\newtcolorbox{tipbox}[1][Tip]{
  colback=tipbg, colframe=primary, fonttitle=\bfseries,
  title=#1, breakable, rounded corners
}
\newtcolorbox{examplebox}[1][Example]{
  colback=examplebg, colframe=primary!70!black, fonttitle=\bfseries,
  title=#1, breakable, rounded corners
}
\newtcolorbox{formulabox}[1][Formula]{
  colback=formulabg, colframe=primary!70!black, fonttitle=\bfseries,
  title=#1, breakable, rounded corners
}
\newtcolorbox{warnbox}[1][Warning]{
  colback=warnbg, colframe=red!60!black, fonttitle=\bfseries,
  title=#1, breakable, rounded corners
}


\title{\color{primary}\textbf{Time \& Work --- Question Pattern Catalog}\\[4pt]\large Every High-Yield Phrasing, Grouped by the Trick That Solves It}
\author{Aptitude Grind}
\date{}

\lhead{\small Time \& Work --- Question Pattern Catalog}
\rhead{\small Aptitude Grind}

\begin{document}
\maketitle
\thispagestyle{fancy}

\noindent Time \& Work questions look varied, but almost every one of them reduces to a handful of
tricks from \texttt{02-fast-tricks-and-shortcuts.tex}. This file catalogues the actual
\textbf{phrasings} you'll meet in TCS NQT (Ninja + Digital) and similar placement tests, grouped by
category, each with the fastest method and a verified mini-example. Skim this before a test to
pattern-match a question to its solution method in seconds.

\section*{Category A --- Two (or More) People Working Together}

\subsection*{A1. Direct Combined-Work Question}
\begin{tipbox}[Phrasing]
``A can finish a work in 6 days, B in 12 days. In how many days can they finish it together?''
\end{tipbox}
\textbf{Method:} $xy/(x+y)$ from \texttt{01-fundamentals.tex} \S3. \\
\textbf{Answer:} $(6 \times 12)/(6+12) = 4$ \textbf{days}.

\subsection*{A2. Three-or-More-Workers Combined Question}
\begin{tipbox}[Phrasing]
``A, B, and C can individually finish a job in 12, 15, and 20 days. Working together, how long will
they take?''
\end{tipbox}
\textbf{Method:} LCM Method (\texttt{02-fast-tricks-and-shortcuts.tex} \S1). LCM(12,15,20) = 60
$\Rightarrow$ rates 5, 4, 3 units/day $\Rightarrow$ combined rate = 12 units/day $\Rightarrow$
\textbf{5 days}.

\section*{Category B --- Efficiency Ratio Problems}

\subsection*{B1. ``X Times as Efficient'' Question}
\begin{tipbox}[Phrasing]
``A is thrice as efficient as B. B alone can complete a work in 24 days. Find the time taken by A
alone, and by A and B together.''
\end{tipbox}
\textbf{Method:} Rate-First Thinking (\S2). B's rate $=1/24$; A's rate $=3/24=1/8$ $\Rightarrow$ A
alone $=$ \textbf{8 days}; together $=(8\times24)/(8+24)=6$ \textbf{days}.

\subsection*{B2. Time-Ratio-to-Efficiency-Ratio Conversion}
\begin{tipbox}[Phrasing]
``The ratio of the time taken by P and Q to complete a task is 4:7. What is the ratio of their
efficiencies?''
\end{tipbox}
\textbf{Method:} Efficiency and time are always inversely proportional
(\texttt{01-fundamentals.tex} \S6) $\Rightarrow$ efficiency ratio $=$ \textbf{7:4}.

\section*{Category C --- Worker Joining or Leaving Midway}

\subsection*{C1. Someone Joins Partway Through}
\begin{tipbox}[Phrasing]
``A can complete a job alone in 18 days. A works alone for 6 days, then B (who alone takes 9 days)
joins. How many total days does the job take?''
\end{tipbox}
\textbf{Method:} Completed-work-then-remaining-work split
(\texttt{02-fast-tricks-and-shortcuts.tex} \S4). A's 6-day share $=6/18=1/3$; remaining $=2/3$;
combined rate $=1/18+1/9=1/6$; extra days needed $=(2/3)/(1/6)=4$. \textbf{Total = 6 + 4 = 10 days.}

\subsection*{C2. Workers Leave Partway Through}
\begin{tipbox}[Phrasing]
``24 workers can finish a job in 18 days. After 6 days, 6 workers leave. In how many more days will
the remaining workers finish the job?''
\end{tipbox}
\textbf{Method:} Worker-Day Method (\texttt{02-fast-tricks-and-shortcuts.tex} \S3). Total
$=24\times18=432$ worker-days; done $=24\times6=144$; remaining $=288$; workers left $=18$
$\Rightarrow$ \textbf{16 more days.}

\section*{Category D --- Men / Women / Children Equivalence}

\subsection*{D1. Given an Equivalence Ratio, Find a Mixed-Group Time}
\begin{tipbox}[Phrasing]
``2 men do the same amount of work as 3 women. 2 men alone can finish a job in 12 days. In how many
days can 4 men and 6 women together finish the same job?''
\end{tipbox}
\textbf{Method:} Convert to a common rate unit first (\texttt{01-fundamentals.tex} \S7). $2\text{
men}=12$ days $\Rightarrow$ 1 man $=1/24$; since $2\text{ men}=3\text{ women}$ in output, 1 woman
$=1/36$. Combined rate of $4\text{ men}+6\text{ women} = 4/24+6/36 = 1/6+1/6 = 1/3 \Rightarrow$
\textbf{3 days.}

\subsection*{D2. Men/Women/Children Together}
\begin{tipbox}[Phrasing]
``If 2 men or 4 women or 6 children can complete a job in the same time, how many days will 1 man, 2
women, and 3 children together take, given 2 men alone take 12 days?''
\end{tipbox}
\textbf{Method:} Same conversion approach as D1 --- every category is first expressed in terms of a
single ``man-unit,'' then rates are added directly.

\section*{Category E --- Pipes and Cisterns}

\subsection*{E1. Fill and Empty Pipe Together}
\begin{tipbox}[Phrasing]
``Pipe A can fill a tank in 8 hours; pipe B can empty it in 12 hours. If both are opened together,
how long will the tank take to fill?''
\end{tipbox}
\textbf{Method:} Net rate (\texttt{01-fundamentals.tex} \S9). $1/8-1/12=1/24 \Rightarrow$
\textbf{24 hours.}

\subsection*{E2. A Pipe Is Closed Partway Through}
\begin{tipbox}[Phrasing]
``Fill pipe A takes 10 hours, outlet pipe B takes 15 hours. Both are opened together, but B is
closed after 3 hours. How long does the tank take to fill in total?''
\end{tipbox}
\textbf{Method:} Two-phase split (\texttt{02-fast-tricks-and-shortcuts.tex} \S8). Net rate while
both open $=1/10-1/15=1/30$; work in 3 hr $=1/10$; remaining $9/10$ filled by A alone at
$1/10$/hr $=$ 9 more hours $\Rightarrow$ \textbf{Total = 12 hours.}

\subsection*{E3. Multiple Fill Pipes, One Outlet}
\begin{tipbox}[Phrasing]
``Three pipes P, Q, R can fill a tank in 12, 15, and 20 minutes respectively, while an outlet pipe S
can empty it in 10 minutes. If all four are opened together, will the tank ever fill?''
\end{tipbox}
\textbf{Method:} Sum every inlet rate, subtract the outlet rate (\texttt{01-fundamentals.tex} \S9).
If the net rate is zero or negative, the tank \textbf{never fills} --- always check the sign of the
net rate before computing a time.

\section*{Category F --- Alternate-Day Working}

\subsection*{F1. Two Workers Alternate Days}
\begin{tipbox}[Phrasing]
``A can complete a job in 10 days, B in 15 days. Working on alternate days, starting with A, in how
many days will the job be completed?''
\end{tipbox}
\textbf{Method:} Cycle-first method (\texttt{02-fast-tricks-and-shortcuts.tex} \S5). 2-day cycle
work $=1/10+1/15=1/6$; since $1/6\times6=1$ exactly, the job finishes in exactly \textbf{12 days} (6
full cycles), with \textbf{no partial final day}.

\subsection*{F2. Alternate Days With a Leftover Fraction}
\begin{tipbox}[Phrasing]
``A can complete a job in 4 days, B in 6 days. Working on alternate days starting with A, how many
full days pass before the job is finished, and who finishes it?''
\end{tipbox}
\textbf{Method:} Same cycle-first method, but here the cycle work ($1/4+1/6=5/12$) does
\textbf{not} divide evenly into 1 --- find how many full cycles fit, then solve the small leftover
fraction with whichever worker's turn comes next (don't assume it always lands on a clean day).

\section*{Category G --- Wages Based on Work}

\subsection*{G1. Split Payment by Work Contributed}
\begin{tipbox}[Phrasing]
``A (10 days alone) and B (15 days alone) work together on a job and are jointly paid Rs.~1500. What is
each person's share?''
\end{tipbox}
\textbf{Method:} Wages-by-work-share (\texttt{02-fast-tricks-and-shortcuts.tex} \S7). Rate ratio
$A:B=3:2 \Rightarrow$ A gets $1500\times3/5=$ Rs.~900, B gets $1500\times2/5=$ Rs.~600.

\subsection*{G2. Three-Way Wage Split With a Helper}
\begin{tipbox}[Phrasing]
``A and B can together do a job in 10 days. With the help of C, they finish it in 6 days. If the
total payment is Rs.~3000, how much should C receive?''
\end{tipbox}
\textbf{Method:} Find C's rate by subtraction (difference-formula logic from
\texttt{01-fundamentals.tex} \S5, applied to a three-person combined rate instead of two), then split
wages by rate ratio exactly as in G1.

\section*{Category H --- Machine / Printer / Robot Work}

\subsection*{H1. Direct Machine-Rate Question}
\begin{tipbox}[Phrasing]
``Printer A can print 1000 pages in 5 hours, Printer B can print the same 1000 pages in 8 hours. If
both print together, how long will it take?''
\end{tipbox}
\textbf{Method:} Identical to a two-worker combined-work question (\texttt{01-fundamentals.tex} \S8)
--- a machine's rate is treated exactly like a person's rate. $(5\times8)/(5+8)=40/13$ hours.

\subsection*{H2. Machine Stopped Partway}
\begin{tipbox}[Phrasing]
``Machine P can complete a task in 8 hours, machine Q in 10 hours, machine R in 12 hours. All three
start together, but P is shut off after 2 hours. In approximately how much more time will Q and R
finish the rest?''
\end{tipbox}
\textbf{Method:} Same two-phase split as E2 --- compute work done by all three in the first phase,
subtract from 1, then solve the remainder using only the machines still running.

\section*{Category I --- Percentage / Fractional Work Completed}

\subsection*{I1. Given a Percentage, Find the Full-Job Time}
\begin{tipbox}[Phrasing]
``A can complete 75\% of a work in 15 days. In how many days can A alone complete the entire work? B
then finishes the remaining 25\% in 4 days --- how many days would B alone take for the full job?''
\end{tipbox}
\textbf{Method:} Scale up directly --- if 75\% takes 15 days, 100\% takes $15/0.75=20$ days for A.
Similarly, B's 25\% in 4 days scales to $4/0.25=16$ days for the full job.

\subsection*{I2. Fractional Work Split Between Two Workers}
\begin{tipbox}[Phrasing]
``A can do 40\% of a job in a certain time; B does the remaining 60\%. If A takes 8 days for their
share, and A and B have the same efficiency, how long does B take for their share?''
\end{tipbox}
\textbf{Method:} Convert percentages to a work ratio ($40:60=2:3$), then scale A's time by that
ratio to find B's time, since equal efficiency means time is proportional to the fraction of work
assigned.

\section*{Category J --- Variable Efficiency (Speed Changes Mid-Job)}

\subsection*{J1. Efficiency Doubles/Triples Partway}
\begin{tipbox}[Phrasing]
``Working at a constant rate, A can finish a job in 10 days. A works at the normal rate for the
first 4 days, then works at double efficiency for the rest. How many total days does the job take?''
\end{tipbox}
\textbf{Method:} Treat this as two separate phases with two different rates, exactly like the
two-phase pipe/machine split in E2/H2. First 4 days at rate $1/10$ complete $2/5$; remaining $3/5$
at double rate ($1/5$/day) takes 3 more days $\Rightarrow$ \textbf{Total = 7 days.}

\subsection*{J2. Efficiency Changes by a Given Percentage}
\begin{tipbox}[Phrasing]
``A normally takes 20 days for a job. If A's efficiency increases by 25\% starting from day one, how
many days will A now take?''
\end{tipbox}
\textbf{Method:} Percentage-to-time formula (\texttt{02-fast-tricks-and-shortcuts.tex} \S6).
$\text{New Time} = 20\times100/125=16$ \textbf{days.}

\section*{Priority Order for Limited Prep Time}
If you're short on time, drill these first --- they cover the large majority of what actually
appears in TCS-style placement tests, roughly in order of frequency:

\begin{formulabox}[Priority List]
\begin{enumerate}[leftmargin=1.4em]
  \item \textbf{A1} --- Two people working together (the base formula everything else builds on)
  \item \textbf{B1} --- Efficiency ratio (``X times as efficient'')
  \item \textbf{C1/C2} --- Someone joins or leaves midway
  \item \textbf{D1} --- Men/Women/Children equivalence
  \item \textbf{E1/E2} --- Pipes and cisterns, including a pipe closing partway
  \item \textbf{F1} --- Alternate-day working
  \item \textbf{G1} --- Wages split by work share
  \item \textbf{H1} --- Machine/Printer/Robot work (same math as A1, different wrapper)
  \item \textbf{I1} --- Percentage of work completed
  \item \textbf{J1/J2} --- Variable efficiency (lower frequency, but appears occasionally in TCS
    Digital)
\end{enumerate}
\end{formulabox}

These ten patterns, combined with the LCM Method and rate-first thinking from
\texttt{02-fast-tricks-and-shortcuts.tex}, cover almost everything TCS typically asks while avoiding
the low-return, highly complex worker puzzles seen in exams like CAT or SSC CGL.

\end{document}
